\section{Đếm phần tử thỏa điều kiện [Cơ bản]}
\subsection{Phân tích \& Thiết kế (M0)}
    \begin{itemize}
    \item \textbf{Input:}
        \begin{itemize}
            \item File \texttt{array.txt} chứa mảng số nguyên: số lượng $n$ và $n$ phần tử.
            \item File \texttt{config.txt} chứa ngưỡng: một số nguyên $k$.
        \end{itemize}

    \item \textbf{Xử lý:}
        \begin{itemize}
            \item \textbf{Case biên:}
            \begin{itemize}
                \item $n = 0 \rightarrow$ mảng rỗng, cả hai biến đếm đều bằng $0$.
                \item Tất cả phần tử trong mảng đều bằng $k \rightarrow$ \texttt{count\_ge\_k} $= n$, \texttt{count\_lt\_k} $= 0$.
            \end{itemize}
            
            \item \textbf{Module 1 (Đọc mảng):}
            \begin{itemize}
                \item Hàm \texttt{readArray(path, vector<int>\&)}
                \item Kiểm tra mở file (\text{fail} $\rightarrow$ \texttt{cerr} + \texttt{return false})
            \end{itemize}

            \item \textbf{Module 2 (Đọc cấu hình k):}
            \begin{itemize}
                \item Hàm \texttt{readConfig(path, int\&)}
                \item Kiểm tra mở file (\text{fail} $\rightarrow$ \texttt{cerr} + \texttt{return false})
            \end{itemize}
            
            \item \textbf{Module 3 (Đếm theo điều kiện):}
            \begin{itemize}
                \item Hàm \texttt{countIf(vector<int>, k, geq)} với \texttt{geq} là biến \texttt{bool} để cờ hiệu đếm lớn hơn hoặc bằng (\texttt{true}) hay nhỏ hơn (\texttt{false}).
            \end{itemize}

            \item \textbf{Module 4 (Ghi output):}
            \begin{itemize}
                \item Hàm \texttt{writeOutput(path, count\_ge\_k, count\_lt\_k)}
            \end{itemize}
        \end{itemize}

        \item \textbf{Output:}
        \begin{lstlisting}
            count_ge_k=<i> / count_lt_k=<i>\end{lstlisting} 
    \end{itemize}

\subsection{Mã nguồn C++11 (Tách module)}
    \begin{lstlisting}
    #include <iostream>
    #include <fstream>
    #include <vector>

    using namespace std;

    // MODULE 1: Doc mang tu file
    bool readArray(string path, vector<int>& a) {
        ifstream file(path);
        if (!file.is_open()) {
            cerr << "Khong mo duoc file " << path << "\n";
            return false;
        }

        int n;
        if (file >> n) {
            int x;
            for (int i = 0; i < n; i++) {
                file >> x;
                a.push_back(x);
            }
        }

        file.close();
        return true;
    }

    // MODULE 2: Doc nguong k tu file config
    bool readConfig(string path, int& k) {
        ifstream file(path);
        if (!file.is_open()) {
            cerr << "Khong mo duoc file " << path << "\n";
            return false;
        }

        file >> k;
        
        file.close();
        return true;
    }

    // MODULE 3: Dem so luong phan tu thoa man
    int countIf(const vector<int>& a, int k, bool geq) {
        int count = 0;
        for (int x : a) {
            if (geq) {
                if (x >= k) count++;
            } else {
                if (x < k) count++;
            }
        }
        return count;
    }

    // MODULE 4: Ghi output
    bool writeOutput(string path, int count_ge_k, int count_lt_k) {
        ofstream file(path);
        if (!file.is_open()) {
            cerr << "Khong ghi duoc file " << path << "\n";
            return false;
        }

        file << "count_ge_k=" << count_ge_k 
            << " / count_lt_k=" << count_lt_k;

        file.close();
        return true;
    }

    // MAIN
    int main() {
        vector<int> a;
        int k;

        if (!readArray("array.txt", a)) {
            return 1;
        }

        if (!readConfig("config.txt", k)) {
            return 1;
        }

        int count_ge = countIf(a, k, true);
        int count_lt = countIf(a, k, false);

        if (!writeOutput("output.txt", count_ge, count_lt)) {
            return 1;
        }

        return 0;
    }
    \end{lstlisting}

\subsection{Kế hoạch kiểm thử (Test Cases)}

    \textbf{Test Case 1 (Thông thường)}

    \textbf{Input (array.txt):}
    \begin{lstlisting}
    6 10 -5 8 15 3 8\end{lstlisting}
    \textbf{Input (config.txt):}
    \begin{lstlisting}
    8\end{lstlisting}
    \textbf{Expected Output (output.txt):}
    \begin{lstlisting}
    count_ge_k=4 / count_lt_k=2\end{lstlisting}

    \vspace{0.5cm}

    \textbf{Test Case 2 (Case biên -- n=0)}

    \textbf{Input (array.txt):}
    \begin{lstlisting}
    0\end{lstlisting}
    \textbf{Input (config.txt):}
    \begin{lstlisting}
    5\end{lstlisting}
    \textbf{Expected Output (output.txt):}
    \begin{lstlisting}
    count_ge_k=0 / count_lt_k=0\end{lstlisting}

    \vspace{0.5cm}
    
    \textbf{Test Case 3 (Case biên -- Tất cả bằng k)}

    \textbf{Input (array.txt):}
    \begin{lstlisting}
    4 7 7 7 7\end{lstlisting}
    \textbf{Input (config.txt):}
    \begin{lstlisting}
    7\end{lstlisting}
    \textbf{Expected Output (output.txt):}
    \begin{lstlisting}
    count_ge_k=4 / count_lt_k=0\end{lstlisting}